%!TEX program = xelatex
\documentclass[12pt, b5paper]{article}
\usepackage{ctex}
\usepackage{exam-zh-question}
\usepackage{exam-zh-choices}
\usepackage{geometry}
\geometry{b5paper,left=2cm,right=1cm,top=2.5cm,bottom=2.5cm}
\usepackage{chemformula} %化学公式宏包
\usepackage[version=4]{mhchem} %化学公式宏包
\usepackage{siunitx} %数字，单位宏包
\sisetup{
	separate-uncertainty = true,
	inter-unit-product = \ensuremath{{}\cdot{}}
} %单位间加小圆点
\usepackage{tikz} %Tikz绘图包
\usetikzlibrary{cd}
\usetikzlibrary{chains}
\usetikzlibrary{arrows}
\usetikzlibrary{arrows.meta} %画箭头
\usetikzlibrary{commutative-diagrams} %交换图
\usetikzlibrary{positioning,automata}
\usepackage[all]{xy}
\usepackage{booktabs} %三线表
\usepackage{multirow} %合并单元格
\usepackage{makecell} %表格内强制换行
\usepackage{array}
\usepackage{booktabs} %控制表格行高
\usepackage{colortbl}  %彩色表格需要加载的宏包
\usepackage{amsmath}
\usepackage{unicode-math}
\usepackage{fontspec} %字体
\newcommand*{\cred}{\textcolor{red}}

\setmainfont{Times New Roman}
\setmathfont{XITS Math}

\setlength{\parindent}{0pt} %无缩进
\usepackage{setspace}
\usepackage{fancyhdr} %设置页码
\usepackage{lastpage} %获取总页数
\pagestyle{fancy} %fancyhdr宏包新增的页面风格
\renewcommand\headrulewidth{0pt} %取消页眉中的装饰分割线
\fancyhf{}
\cfoot{\footnotesize 第 \thepage\ 页(共 \pageref{LastPage}页)}%当前页 of 总页数





\begin{document}
\begin{spacing}{1.625}

\centerline{{\huge 河北农业大学课程考试试卷}}

\centerline{\large 2022-2023学年第\underline{\ 1\ }学期}

考试科目：\underline{物理化学与胶体化学} \quad 姓名：\underline{ \qquad \qquad} \quad 学号：\underline{ \qquad \qquad \qquad}

学院：\underline{理工学院} \quad 专业班级：\underline{\qquad \qquad \qquad \qquad} \quad 卷别：\cred{A} 考核方式：\cred{开卷}

\bigskip


1. 简答题(15分)

物理化学和生活息息相关。请你举两个例子，说明物理化学在生产生活中的重要意义。

\bigskip

\bigskip

2. 简答题(15分)

化学系流传着一句话，叫做“有机物化，必有一挂”。无机、有机、分析等学科，它们都有特殊的研究对象，而物化则着重于研究更具普遍性的化学运动的本质特征和内在规律性。你是如何学习物理化学的？你认为怎样才能学好物理化学？

\bigskip

\bigskip

3. 简答题(10分)

对于某一相平衡体系的物种数与独立组分数，从不同的角度出发考虑问题，体系的物种数可以不同，但独立组分数一定是一样的。试举例说明。

\bigskip

\bigskip

4. 主观题(20分)

乙烷脱氢生产乙烯 \ch{C2H6 (g) -> C2H4 (g) + H2 (g)}

(1) 判断该反应在常温常压下能否自发进行。

(2) 判断该反应在500K下能否自发进行。相关热力学数据查课本附录。

\bigskip

\bigskip


5. 主观题(20分)

试计算在珠穆朗玛峰顶上的沸水中煮熟一个鸡蛋需要多长时间。

己知组成蛋的蛋白质的热变作用为一级反应，其活化能约为\SI{85}{kJ.mol^{-1}}，在$p^\ominus$下的沸水中煮熟一个鸡蛋需要\SI{10}{min}。水的摩尔蒸发热$\Delta H_\mathrm{m} = \SI{41.05}{kJ.mol^{-1}}$。珠穆朗玛峰顶大气压力$p=\SI{31}{kPa}$。

\bigskip

\bigskip

6. 主观题(20分)

请你设计一道综合计算题，要求包含化学热力学和化学动力学的知识点。将答案一并附上。


\end{spacing}

% \begin{question}
% 	若过点 $(a, b)$ 可作曲线 $y = ^x$ 的两条切线，则 \paren
% 	\begin{choices}
% 	  \item $^b < a$
% 	  \item $^a < b$
% 	  \item $0 < a < ^b$
% 	  \item $0 < b < ^a$
% 	\end{choices}
%   \end{question}

%   \begin{problem}[points = 12]
% 	已知函数 $f(x) = x (1 - \ln x)$。
% 	\begin{enumerate}
% 	  \item 讨论 $f(x)$ 的单调性；
% 	  \item 设 $a$，$b$ 为两个不相等的正数，且 $b \ln a - a \ln b = a - b$，
% 		证明：$2 < \frac{1}{a} + \frac{1}{b} < $。
% 	\end{enumerate}
%   \end{problem}
  
% 	函数的定义域为 $(0, +\infty)$,

% 	又 $f^{\prime}(x) = 1 - \ln x-1 = -\ln x$, \score{2}

% 	当 $x \in(0, 1)$ 时, $f^{\prime}(x) > 0$, 当 $x \in(1, +\infty)$ 时, $f^{\prime}(x) < 0$,
% 	故 $f(x)$ 的递增区间为 $(0,1)$, 递减区间为 $(1, +\infty)$.

% \begin{question}[index=1]
%   已知函数 $f(x) = x^3 (a \cdot 2^x - 2^{-x})$ 是偶函数，则 $a = $ \fillin[$1$] 。
% \end{question}




\end{document}
